%! Logarithmic Function Manipulation — Unit 2, Lesson 2.12 — Slides (Beamer)
\documentclass[aspectratio=169, 11pt]{beamer}
\usepackage{precalculus-beamer}   % provides \navyheader{} and \sectionlabel[color]{}

\begin{document}

% TITLE SLIDE
{
\setbeamercolor{background canvas}{bg=navy}
\begin{frame}[plain]
  \vspace{2.8cm}\hspace{0.55cm}%
  \begin{minipage}{0.9\textwidth}
    {\color{goldacc}\bfseries\small LESSON 2.12}\\[0.5cm]
    {\color{white}\bfseries\LARGE Logarithmic Function Manipulation}\\[1.2cm]
    {\color{skymid}\small Precalculus~$\cdot$~Unit 2: Exponential and Logarithmic Functions}
  \end{minipage}
\end{frame}
}

% SLIDE 1 — Hook
\begin{frame}[t, plain]
  \navyheader{Hook: Can You Predict the Pattern?}

  \vspace{0.3cm}
  We know: \quad $\log_2(8) = 3$ \quad and \quad $\log_2(4) = 2$

  \vspace{0.5cm}
  \begin{block}{Think--Pair--Share}
    What do you predict $\log_2(8 \cdot 4) = \log_2(32)$ equals?\\[6pt]
    Can you write it as a sum of the two values above?
  \end{block}

  \vspace{0.5cm}
  \textbf{Verify:} $\log_2(32) = ?$ because $2^{\,?} = 32$

  \vspace{0.4cm}
  \textit{Today we generalize: when can we break a log of a product into a sum?}
\end{frame}

% SLIDE 2 — Product Property
\begin{frame}[t, plain]
  \navyheader{Product Property (EK 2.12.A.1)}

  \vspace{0.2cm}
  \begin{block}{Product Property}
    \[ \log_b(xy) = \log_b(x) + \log_b(y) \]
  \end{block}

  \vspace{0.3cm}
  \textbf{Examples:}
  \begin{itemize}
    \setlength{\itemsep}{6pt}
    \item $\log_2(8 \cdot 4) = \log_2(8) + \log_2(4) = 3 + 2 = 5$
    \item $\log_3(9 \cdot 27) = \log_3(9) + \log_3(27) = 2 + 3 = 5$
  \end{itemize}

  \vspace{0.35cm}
  \begin{block}{Graphical Connection}
    $f(x) = \log_b(kx) = \underbrace{\log_b(k)}_{a} + \log_b(x)$\\[4pt]
    A \textbf{horizontal dilation} by factor $k$ is equivalent to a
    \textbf{vertical translation} up by $\log_b(k)$.
  \end{block}
\end{frame}

% SLIDE 3 — Power Property
\begin{frame}[t, plain]
  \navyheader{Power Property (EK 2.12.A.2)}

  \vspace{0.2cm}
  \begin{block}{Power Property}
    \[ \log_b(x^n) = n \cdot \log_b(x) \]
  \end{block}

  \vspace{0.3cm}
  \textbf{Examples:}
  \begin{itemize}
    \setlength{\itemsep}{6pt}
    \item $\log_3(9^2) = 2 \cdot \log_3(9) = 2 \cdot 2 = 4$
    \item $\log_2(x^4) = 4\log_2(x)$
  \end{itemize}

  \vspace{0.35cm}
  \begin{block}{Graphical Connection}
    $f(x) = \log_b(x^k) = k\log_b(x)$\\[4pt]
    Raising the input to a power $k$ is equivalent to a
    \textbf{vertical dilation} by factor $k$.
  \end{block}
\end{frame}

% SLIDE 4 — Change of Base
\begin{frame}[t, plain]
  \navyheader{Change of Base (EK 2.12.A.3)}

  \vspace{0.2cm}
  \begin{block}{Change-of-Base Formula}
    \[ \log_b(x) = \frac{\log_a(x)}{\log_a(b)} \quad \text{for any valid base } a \]
  \end{block}

  \vspace{0.3cm}
  \textbf{Calculator form} (using common log):
  \[ \log_5(17) = \frac{\log(17)}{\log(5)} \approx 1.760 \]

  \vspace{0.25cm}
  \begin{block}{Big Implication}
    $\log_b(x) = \dfrac{1}{\log_{10}(b)} \cdot \log_{10}(x)$ \quad --- a vertical dilation of $\log_{10}(x)$\\[4pt]
    \textbf{All logarithmic functions are vertical dilations of each other.}
  \end{block}
\end{frame}

% SLIDE 5 — Natural Logarithm
\begin{frame}[t, plain]
  \navyheader{Natural Logarithm (EK 2.12.A.4)}

  \vspace{0.3cm}
  \begin{block}{Definition}
    $\ln(x) = \log_e(x)$, where $e \approx 2.718\ldots$ (Euler's number)
  \end{block}

  \vspace{0.4cm}
  \textbf{Key evaluations} (no calculator):
  \begin{itemize}
    \setlength{\itemsep}{8pt}
    \item $\ln(e) = \log_e(e) = 1$ \quad (since $e^1 = e$)
    \item $\ln(e^3) = 3\ln(e) = 3$ \quad (power property)
    \item $\ln(1) = 0$ \quad (since $e^0 = 1$)
    \item $\ln(e^{-2}) = -2$ \quad (power property)
  \end{itemize}

  \vspace{0.3cm}
  \textit{All product and power properties apply to $\ln$ exactly as to any $\log_b$.}
\end{frame}

% SLIDE 6 — Class Practice
\begin{frame}[t, plain]
  \navyheader{Class Practice}

  \vspace{0.2cm}
  Try these with a partner. Show each step.

  \vspace{0.3cm}
  \begin{enumerate}
    \setlength{\itemsep}{10pt}
    \item Rewrite $\log_2(8x^4)$ using the product and power properties.
    \item Rewrite $\log_4(16x^2)$ in the form $a + b\log_4(x)$. Find $a$ and $b$.
    \item Evaluate $\log_7(50)$ using change of base. Express as a ratio, then find the decimal.
    \item Evaluate $\ln(e^5)$, $\ln(1)$, and $\ln(e^{-2})$ without a calculator.
  \end{enumerate}

  \vspace{0.4cm}
  \begin{block}{Exit Ticket}
    Questions 2 and 3 are your exit ticket today.
  \end{block}
\end{frame}

\end{document}
